\begin{table}[h]
    \caption{Classification of feature engineering techniques used in the reviewed studies.}
    \label{tab:PS_RQ_3_Feature_engineering_classification}
    \centering
    %{\fontfamily{phv}\selectfont   % inizio gruppo → Helvetica/Arial-like
    \scriptsize                 % o \footnotesize, \small…
    \resizebox{0.9\textwidth}{!}{
    \begin{tabular}{|
  >{\centering\arraybackslash}m{5cm}
  |>{\centering\arraybackslash}m{8cm}|}
    \hline
    \rowcolor[HTML]{D5E8D4} 
    \rule[-1.5ex]{0pt}{4ex}\textbf{Feature Engineering Method }
  & \rule[-1.5ex]{0pt}{4ex}\textbf{Study IDs} \\ \hline
    RFM and its extensions & S3, S5, S6, S9, S8, S17,
S21, S22, S28, S30, S31,
S33, S35, S36, S37, S41 \\ \hline
    Feature selection & S2, S4, S9, S12, S18, S21,
S19, S23, S27, S32, S35,
S38, S40 \\ \hline
    Feature transformation/
aggregation & S5, S9, S11, S15, S16,
S22, S23, S33, S35, S37 \\ \hline
    Customer segmentation
techniques & S19, S24, S27, S29, S30,
S37 \\ \hline
    Dimensionality reduction  & S9, S16, S19, S22, S23 \\ \hline
    \end{tabular}
    }
%  }% fine gruppo → torna al font normale qui    
\end{table}